%================================
% Vocabulary & Language Focus — XB2 (Unit 3)
% 编译方式: xelatex main.tex (需运行两遍以生成目录)
%================================
\documentclass[../main.tex]{subfiles}

\begin{document}

\section{Unit 3}

\vw{specialize}{/ˈspeʃəlaɪz/}{vi.}

\sense{vi.}{专门研究（或从事），专攻}
\col{specialize in sth}
\ex{The doctor specializes in heart surgery.}
\ex{The shop specializes in handmade shoes.}
\col{specialist (n.) 专家}
\ex{He is a specialist in ancient history.}
\col{specialized (adj.) 专门的}
\ex{specialized equipment}

\begin{origin}
来自拉丁语 \textit{specialis}（特别的），源于 \textit{species}（种类、外观）。同源词：\textbf{species}（物种）、\textbf{special}、\textbf{specific}。"在某一类里做到专精"就是 specialize。
\end{origin}

\vw{being}{/ˈbiːɪŋ/}{n.}

\sense{n.}{生物}
\col{a human being 人}
\ex{No living being can survive without water.}
\col{intelligent beings 有智慧的生物}

\sense{n.}{存在}
\col{come into being 开始存在，产生}
\ex{The universe came into being billions of years ago.}
\col{bring sth into being 使产生}

\sense{n.}{身心}
\col{the whole of one's being 全身心}
\ex{She threw her whole being into the performance.}

\col{well-being 幸福，安康}
\ex{Regular exercise improves your sense of well-being.}

\begin{origin}
being 是动词 be 的现在分词名词化，译自希腊哲学的"存在"（\textit{to on}）——亚里士多德谓"存在之为存在"是形而上学的研究对象。英语语法中 being 还构成进行时被动式（\textit{being built}）与动名词（\textit{without being noticed}）。
\end{origin}

\vw{lifespan}{/ˈlaɪfspæn/}{n.}

\sense{n.}{寿命，有效期}
\col{a lifespan of + 时长}
\ex{The average lifespan of a cat is about 15 years.}
\col{the lifespan of a product 产品有效期}
\ex{The lifespan of these batteries is two years.}

\begin{origin}
life + span。span 源自古英语 \textit{spann}——"手掌张开的一拃"，由"一拃的长度"引申为"一段时间"。同源词：\textbf{span}（跨度；持续）。"生命的一拃"就是寿命。
\end{origin}

\vw{affect}{/əˈfekt/}{vt.}

\sense{vt.}{影响}
\col{affect sb/sth}
\ex{The new law will affect millions of families.}
\col{be affected by sth}
\ex{Tourism in the region was badly affected by the earthquake.}

\sense{vt.}{侵袭，使感染}
\ex{The disease affects the nervous system.}

\sense{vt.}{深深打动}
\col{be deeply/profoundly affected}
\ex{We were deeply affected by her story.}

\begin{lex}{affect / effect}
\textbf{affect} 是动词（影响）：\textit{The rain affected the match.}；\textbf{effect} 是名词（影响、效果）：\textit{have an effect on sth}（对……有影响）。拼写只差一字母，是最经典的易混词对。注意 effect 也可作动词，但表"实现、引起"（\textit{effect changes} 实现变革），属正式用语。
\end{lex}

\begin{origin}
来自拉丁语 \textit{afficere}（施加作用、影响）：\textit{ad-}（向）+ \textit{facere}（做）。同源词：\textbf{effect}（e- 出 + facere $\to$ 做出来的结果）、\textbf{fact}、\textbf{factory}。"对……施加作用"即影响。
\end{origin}

\vw{contract}{/kənˈtrækt/ \ vt. \ \ /ˈkɒntrækt/ \ n.}{vt. \& n.}

\sense{vt. /kənˈtrækt/}{感染（疾病）}
\col{contract a disease/virus}
\ex{He contracted malaria while travelling in Africa.}

\sense{vt.}{与……订立合同（或契约）}
\col{contract with sb to do sth}
\ex{The firm contracted with a local builder to construct the warehouse.}

\sense{n. /ˈkɒntrækt/}{合同，合约}
\col{sign/draw up a contract}
\ex{She signed a two-year contract with the club.}
\col{be under contract 在合同期内}
\col{breach of contract 违约}
\ex{The company was sued for breach of contract.}

\begin{gram}[名动异读：重音转移构词法]
\textit{contract} 名词重音在前（/\textbf{ˈkɒn}trækt/ 合同），动词重音在后（/kən\textbf{ˈtrækt}/ 感染；缔约）。英语中一批双音节词遵循"动词重音在后、名词重音在前"的规律：\textit{re\textbf{cord}}（v. 记录）/ \textbf{re}cord（n. 唱片）、\textit{in\textbf{crease}}（v. 增加）/ \textbf{in}crease（n. 增加）、\textit{de\textbf{crease}}（v. 减少）/ \textbf{de}crease（n. 减少）。判断词性可先看重音位置。
\end{gram}

\begin{origin}
来自拉丁语 \textit{contrahere}（拉到一起）：\textit{con-}（一起）+ \textit{trahere}（拉）。同源词：\textbf{tractor}（拖拉机：拉东西的机器）、\textbf{attract}（吸引）、\textbf{abstract}（抽出）、\textbf{extract}（提取）。三个义项共享"拉"的意象：缔约（双方拉拢）、感染（把病拉上身）、收缩（肌肉拉紧）。
\end{origin}

\vw{cancer}{/ˈkænsə(r)/}{n.}

\sense{n.}{癌症}
\col{lung/breast cancer}
\ex{He died of lung cancer.}
\col{cancer cells 癌细胞}

\sense{n.}{邪恶，毒瘤}
\ex{Corruption is a cancer in the body politic.}

\begin{origin}
来自拉丁语 \textit{cancer}（螃蟹），对应希腊语 \textit{karkinos}。古希腊名医希波克拉底观察恶性肿瘤，见肿块周围血管怒张如蟹爪，遂以"螃蟹"名之。星座"巨蟹座"（Cancer）同源，词根还有 \textbf{carcinoma}（癌）。
\end{origin}

\vw{identify}{/aɪˈdentɪfaɪ/}{vt.}

\sense{vt.}{找到，发现}
\col{identify the cause/problem}
\ex{Researchers identified the cause of the outbreak.}

\sense{vt.}{确认，认出}
\col{identify sb/sth as sth}
\ex{The victim was identified as a 30-year-old man.}

\sense{vt.}{显示，说明身份}
\ex{The police asked the man to identify himself.}

\col{identify with sb 认同，与……产生共鸣}
\ex{Many young people identify with the hero of the novel.}

\begin{origin}
来自拉丁语 \textit{idem}（同一个）$\to$ \textit{identitas}（同一性）$\to$ identify。"确认是同一个人/物"即辨认。同源词：\textbf{identical}（完全相同的）、\textbf{identity}（身份）。
\end{origin}

\vw{relate}{/rɪˈleɪt/}{vi. \& vt.}

\sense{vt.}{联系，使有联系}
\col{relate A to B}
\ex{It's hard to relate the two events.}

\sense{vt.}{叙述，讲述}
\col{relate a story/experience}
\ex{He related the whole story to us.}

\sense{vi.}{relate to 涉及，与……相关}
\ex{The new policy relates to all employees.}

\sense{vi.}{relate to sb 理解，产生共鸣}
\ex{Many students find it hard to relate to the novel.}

\col{related (adj.) 相关的}
\ex{related problems}

\begin{origin}
来自拉丁语 \textit{relatus}（带回来的），是 \textit{referre}（refer）的过去分词：\textit{re-}（回）+ \textit{ferre}（携带）。同源词：\textbf{refer}、\textbf{relative}（相关的；亲戚）、\textbf{reference}。"把……带回来"——讲述往事、把两件事带在一起，皆由此。
\end{origin}

\vw{means}{/miːnz/}{n.}

\sense{n.}{(pl. means) 方法，方式}
\col{a means of doing sth 做……的方法}
\ex{Email is a quick means of communication.}
\col{by means of 借助，通过}
\ex{The message was sent by means of a courier.}
\col{by all means 当然可以}
\ex{"May I sit here?" "By all means!"}
\col{by no means 决不，绝不}
\ex{This is by no means the first time the problem has appeared.}

\sense{n.}{财富，钱财}
\col{a man of means 有钱人}
\ex{He is a man of means.}
\col{beyond one's means 超出某人的财力}
\ex{He was living beyond his means.}

\begin{gram}[means 单复数同形]
\textit{means} 单复数同形：\textit{a means} / \textit{two means}（\textcolor{WildStrawberry!80!black}{\ding{55}} \textit{a mean}）。作"方法"讲时可数；作"钱财"讲不可数。固定短语 \textit{by all means}、\textit{by no means}、\textit{by means of} 是高频考点，注意 \textit{by no means} 置于句首时句子部分倒装：\textit{By no means is this the only solution.}
\end{gram}

\begin{origin}
来自 mean（中间）——古法语 \textit{meien}，源于拉丁语 \textit{medius}（中间的）。"取中间之道"就是方法。同源词：\textbf{medium}（媒介）、\textbf{median}（中位数）、\textbf{middle}。
\end{origin}

\vw{restore}{/rɪˈstɔː(r)/}{vt.}

\sense{vt.}{恢复}
\col{restore peace/order/confidence}
\ex{The army restored order in the city.}
\ex{Rest is the best way to restore your health.}

\sense{vt.}{修复}
\col{restore sth to its former state}
\ex{The church has been restored to its original condition.}

\sense{vt.}{使复原，使复位}
\ex{The police restored the stolen painting to its owner.}

\begin{origin}
来自拉丁语 \textit{restaurare}（重建）：\textit{re-}（再）+ \textit{instaurare}（建立）。同源词：\textbf{restaurant}（餐馆）——法语 \textit{restaurer}（恢复体力），餐馆最初是"让人恢复元气的地方"。restore 与 restaurant 同源。
\end{origin}

\vw{ultimately}{/ˈʌltɪmətli/}{adv.}

\sense{adv.}{最终；根本上}
\ex{Ultimately, the decision is yours.}
\ex{Ultimately, success depends on hard work.}
\col{ultimate (adj.) 最终的；极致的}
\ex{the ultimate goal 最终目标}

\begin{origin}
来自拉丁语 \textit{ultimus}（最后的）——\textit{ulter}（那边的）的最高级。同源词：\textbf{ultimate}、\textbf{ultra-}（超，如 ultraviolet 紫外线）、\textbf{ulterior}（更远的；不可告人的）。“走到最远的那一端”就是最终。
\end{origin}

\vw{fundamental}{/ˌfʌndəˈmentl/}{adj. \& n.}

\sense{adj.}{十分重大的，根本的；基础的，基本的}
\col{a fundamental change/difference}
\ex{There is a fundamental difference between the two systems.}
\col{be fundamental to sth 对……至关重要}
\ex{Trust is fundamental to any relationship.}

\sense{n.}{基本规律，根本法则，基础}
\col{the fundamentals of sth}
\ex{the fundamentals of physics}

\begin{origin}
来自拉丁语 \textit{fundamentum}（地基），源于 \textit{fundare}（打地基），词根 \textit{fundus}（底部）。同源词：\textbf{foundation}（地基；基金会）、\textbf{fund}（基金）、\textbf{found}（创建）。“地基上的东西”就是根本的。
\end{origin}

\vw{pace}{/peɪs/}{n. \& vi. \& vt.}

\sense{n.}{发生的速度，节奏}
\col{at a fast/slow pace}
\ex{The project is progressing at a slow pace.}
\col{keep pace with sth 跟上}
\ex{Wages have not kept pace with inflation.}
\col{at one's own pace 按自己的节奏}
\ex{Learn at your own pace.}

\sense{n.}{步速}
\col{quicken one's pace 加快脚步}
\ex{He quickened his pace as it began to rain.}

\sense{n.}{快节奏}
\col{the pace of life 生活节奏}
\ex{Many people dislike the fast pace of city life.}

\sense{vi.}{来回踱步，走来走去}
\col{pace up and down}
\ex{He paced up and down the room, waiting for news.}

\sense{vt.}{确定速度，调整节奏}
\col{pace oneself 调整自己的节奏}
\ex{She paced herself carefully over the 10-kilometre run.}

\begin{origin}
来自拉丁语 \textit{passus}（一步），源于 \textit{pandere}（张开）——"迈开的一步"。同源词：\textbf{pass}（经过）、\textbf{passage}（通道）、\textbf{compass}（圆规：用步伐丈量）。“一步”引申出“步速、节奏”。
\end{origin}

\vw{arise}{/əˈraɪz/}{vi.}

\sense{vi.}{(arose, arisen) 发生，出现}
\ex{A new difficulty has arisen.}
\ex{Questions arise as more data becomes available.}

\sense{vi.}{（由……）引起，（因……）产生}
\col{arise from/out of sth}
\ex{Many problems arise from misunderstandings.}

\begin{lex}{arise / rise / raise / arouse}
\textbf{arise}（arose, arisen）不及物，表问题、情况"出现、产生"：\textit{problems arise}；\textbf{rise}（rose, risen）不及物，表"上升"：\textit{the sun rises}、\textit{prices rise}；\textbf{raise}（raised, raised）及物，表"举起、提高、抚养"：\textit{raise your hand}、\textit{raise prices}；\textbf{arouse}（aroused, aroused）及物，表"唤起"：\textit{arouse suspicion}。四词中只有 arise 与 arouse 以 a- 开头，常被混淆。
\end{lex}

\begin{origin}
源自古英语 \textit{arīsan}（起身）：\textit{a-}（向上）+ \textit{rīsan}（起来）。同源词：\textbf{rise}、\textbf{raise}、\textbf{rouse}（唤醒）。日耳曼语词根意为"站起"。
\end{origin}

\vw{restrict}{/rɪˈstrɪkt/}{vt.}

\sense{vt.}{限制，限定}
\col{restrict sth to sth}
\ex{The school restricts mobile phone use to certain areas.}
\col{be restricted to 仅限于}
\ex{Entry is restricted to members only.}

\sense{vt.}{束缚，妨碍；约束，管束}
\col{restrict sb from doing sth}
\ex{The new rules restrict drivers from using phones.}

\col{restriction (n.) 限制}
\ex{impose restrictions on sth}

\begin{lex}{restrict / limit / confine}
\textbf{restrict} 侧重"限定范围"（\textit{restrict access} 限制进入）；\textbf{limit} 侧重"设上限"（\textit{limit your answer to 200 words}）；\textbf{confine} 侧重"限制在空间/活动范围内"，常含"禁锢"义（\textit{confined to bed} 卧床不起）。
\end{lex}

\begin{origin}
来自拉丁语 \textit{restringere}（捆紧）：\textit{re-}（回）+ \textit{stringere}（捆、拉紧）。同源词：\textbf{strict}（严格的）、\textbf{restrain}（抑制）、\textbf{constraint}（约束）、\textbf{stringent}（严苛的）。"捆起来"就是限制。
\end{origin}

\vw{prohibit}{/prəˈhɪbɪt/}{vt.}

\sense{vt.}{（尤指以法令）禁止}
\col{prohibit sth}
\ex{Smoking is prohibited in public buildings.}

\sense{vt.}{使不可能}
\ex{The steep terrain prohibited any further advance.}

\col{prohibit sb from doing sth 禁止某人做某事}
\ex{The law prohibits drivers from using mobile phones.}

\begin{lex}{prohibit / ban / forbid}
\textbf{prohibit} 最正式，常指法律禁止（\textit{prohibited by law}）；\textbf{ban} 常见，指官方明文禁止（\textit{ban smoking}）；\textbf{forbid} 侧重个人/权威的命令（\textit{I forbid you to go}）。句型：\textit{prohibit/ban sb from doing sth}；\textit{forbid sb to do sth}。
\end{lex}

\begin{origin}
来自拉丁语 \textit{prohibere}（挡在前面）：\textit{pro-}（在前）+ \textit{habere}（握住）。"伸手挡住去路"就是禁止。同源词：\textbf{inhibit}（抑制）、\textbf{exhibit}（展出：拿出来给人看）、\textbf{habit}（习惯：握住的状态）。
\end{origin}

\vw{resistance}{/rɪˈzɪstəns/}{n.}

\sense{n.}{抵抗力}
\col{resistance to sth}
\ex{The body builds up resistance to the virus.}

\sense{n.}{抵制，反对}
\col{resistance to sth}
\ex{There was strong resistance to the new tax.}

\sense{n.}{抵抗，反抗}
\col{put up resistance}
\ex{The defenders put up fierce resistance.}

\col{resist (v.) 抵抗；忍住}
\ex{I couldn't resist the chocolate cake.}

\begin{origin}
来自拉丁语 \textit{resistere}（站在对面）：\textit{re-}（相对）+ \textit{sistere}（站立）。同源词：\textbf{resist}、\textbf{insist}（坚持）、\textbf{consist}（组成）。“站在对面不动”就是抵抗。
\end{origin}

\vw{flu}{/fluː/}{n.}

\sense{n.}{流行性感冒，流感}
\col{have/get the flu 得流感}
\ex{She's in bed with the flu.}
\col{a flu jab 流感疫苗注射}
\ex{He got a flu jab before winter.}

\begin{origin}
flu 是 \textit{influenza} 的缩略。influenza 借自意大利语，源于拉丁语 \textit{influentia}（影响），词根 \textit{influere}（流入）。中世纪星相学认为星体"流入"的气息影响人间，致人得病——18 世纪流感大流行时，意大利人用"星象影响"命名此病。
\end{origin}

\vw{characteristic}{/ˌkærəktəˈrɪstɪk/}{n. \& adj.}

\sense{n.}{特征，特点}
\col{a characteristic of sth}
\ex{Politeness is a characteristic of the local people.}

\sense{adj.}{典型的，独特的}
\col{be characteristic of sb/sth}
\ex{That smile is characteristic of her.}
\ex{the characteristic taste of olive oil}

\begin{origin}
来自希腊语 \textit{kharaktēr}（雕刻工具、刻痕），源于 \textit{kharassein}（雕刻）。同源词：\textbf{character}（性格；汉字）、\textbf{characterize}（刻画）。"刻下的印记"就是特征——性格即时间刻下的痕迹。
\end{origin}

\vw{wrestle}{/ˈresl/}{vi. \& vt.}

\sense{vi. \& vt.}{奋力对付，努力处理}
\col{wrestle with a problem/dilemma}
\ex{He wrestled with the problem for days.}
\ex{The country is wrestling with high unemployment.}

\sense{vi. \& vt.}{摔跤}
\col{wrestle with sb}
\ex{The boys wrestled on the grass.}

\col{wrestle with 努力处理，全力解决}
\ex{We are still wrestling with the question of funding.}

\begin{origin}
源自古英语 \textit{wræstlian}，源于 \textit{wræstan}（扭），词根可追溯到原始印欧语 *\textit{wreit-}（扭转）。同源词：\textbf{wrist}（手腕——能扭动的关节）、\textbf{writhe}（扭动）。"扭打"就是摔跤，"扭住难题不放"就是奋力对付。
\end{origin}

\vp{put sth in place}

\sense{v.}{准备就绪；落实到位}
\col{put measures/rules/plans in place}
\ex{New safety measures have been put in place.}
\ex{The government put emergency plans in place before the storm.}
\col{in place 就位，到位}
\ex{Everything is in place for the ceremony.}

\begin{origin}
in place 字面义"在应有的位置上"。put sth in place 即"把……放到该放的位置上"，引申为"落实（措施、制度）"。
\end{origin}

\vw{debate}{/dɪˈbeɪt/}{vt. \& vi. \& n.}

\sense{vt. \& vi.}{讨论，辩论}
\col{debate a topic/proposal}
\ex{The committee debated the proposal for hours.}

\sense{vt. \& vi.}{仔细考虑，思考}
\col{debate whether/what/how to do sth}
\ex{She debated whether to accept the offer.}

\sense{n.}{讨论，争论；辩论}
\col{under debate 在讨论中}
\ex{The issue is still under debate.}
\col{a heated/lively debate}
\ex{The plan sparked a heated debate.}

\begin{origin}
来自拉丁语 \textit{battuere}（打）$\to$ 古法语 \textit{debatre}（打斗）。同源词：\textbf{batter}（连续击打）、\textbf{battle}（战斗）。"用言语打斗"就是辩论。
\end{origin}

\vp{without doubt}

\sense{adv.}{毫无疑问，的确}
\ex{She is without doubt the best player on the team.}
\col{no doubt 无疑地}
\ex{No doubt he will be late again.}
\col{beyond (a) doubt 确凿无疑}
\ex{His guilt is beyond doubt.}

\begin{origin}
doubt 来自拉丁语 \textit{dubitare}（犹豫），源于 \textit{dubius}（摇摆不定的）——词根 duo（二），“在两者之间摇摆”。同源词：\textbf{double}（双倍）、\textbf{dual}（双重的）。"在是与非之间摇摆"就是怀疑。
\end{origin}

\vw{ripe}{/raɪp/}{adj.}

\sense{adj.}{成熟的}
\col{ripe fruit/cheese}
\ex{These peaches are ripe and sweet.}
\col{a ripe old age 高龄}
\ex{She lived to a ripe old age of 95.}

\sense{adj.}{时机成熟的}
\col{the time is ripe for sth}
\ex{The time is ripe for reform.}

\begin{origin}
源自古英语 \textit{rīpe}，与 \textit{rīpan}（收割）同根——"可以收割了"就是成熟。同源词：\textbf{reap}（收割）。ripe 与 reap 同源。
\end{origin}

\vw{withdraw}{/wɪðˈdrɔː/}{vi. \& vt.}

\sense{vi. \& vt.}{(withdrew, withdrawn) 脱离（社会），不与人交往}
\col{withdraw into oneself 变得孤僻}
\ex{After the trauma, he withdrew into himself.}

\sense{vt.}{（使）撤回；（使）退出}
\col{withdraw from sth}
\ex{He withdrew from the race due to injury.}
\col{withdraw a statement/claim 撤回声明}
\ex{The company withdrew its claim.}

\sense{vt.}{停止提供}
\ex{The government withdrew funding from the programme.}

\col{withdraw money 取钱}
\ex{She withdrew 200 dollars from the bank.}
\col{withdrawal (n.) 撤回；取款}

\begin{origin}
中古英语 \textit{withdrawen}：\textit{with}（离开、往回）+ \textit{drawen}（拉）——with- 是表"离开"的古老前缀（同 \textbf{withstand} 站在一旁 $\to$ 抵抗）。"往回拉"就是撤回、退出。
\end{origin}

\vw{equip}{/ɪˈkwɪp/}{vt.}

\sense{vt.}{配备}
\col{be equipped with sth}
\ex{The laboratory is equipped with modern instruments.}

\sense{vt.}{使有能力}
\col{equip sb with sth}
\ex{The course equips students with practical skills.}

\col{equipment (n.) 设备（不可数）}
\ex{a piece of equipment}

\begin{gram}[equipment 不可数]
\textit{equipment} 是不可数名词：\textcolor{WildStrawberry!80!black}{\ding{55}} \textit{an equipment} / \textit{equipments}；应说 \textit{a piece of equipment} 或 \textit{some equipment}。同类不可数名词：\textbf{furniture}、\textbf{luggage}、\textbf{information}。
\end{gram}

\begin{origin}
来自古法语 \textit{esquiper}（装备船只），源于古诺斯语 \textit{skipa}（布置、安排），词根 \textit{skip}（船）。同源词：\textbf{ship}（船）、\textbf{skipper}（船长）。"出海前的布置"就是装备——equip 与 ship 同源。
\end{origin}

\vw{artificial}{/ˌɑːtɪˈfɪʃl/}{adj.}

\sense{adj.}{人工的，人造的}
\col{artificial intelligence 人工智能}
\col{artificial flowers/limbs 假花；假肢}
\ex{The lake is artificial, built in the 1980s.}

\sense{adj.}{人为的}
\col{an artificial price 人为定价}

\sense{adj.}{虚假的}
\col{an artificial smile 假笑}
\ex{She gave an artificial smile and walked away.}

\begin{origin}
来自拉丁语 \textit{artificialis}，源于 \textit{artificium}（手艺、技艺），词根 \textit{artifex}（工匠）：\textit{ars}（艺术、手艺）+ \textit{facere}（做）。同源词：\textbf{art}、\textbf{artist}、\textbf{artifact}（人工制品）。"靠手艺做出来的"就是人造的。
\end{origin}

\vw{limb}{/lɪm/}{n.}

\sense{n.}{肢，臂，腿}
\col{artificial limbs 假肢}
\ex{He lost a limb in the accident.}

\col{out on a limb 处境孤立；冒险}
\ex{He went out on a limb to defend the policy.}

\col{limb 中的 b 不发音，同 comb、dumb}

\begin{origin}
源自古英语 \textit{lim}（肢），日耳曼语词根，确切来源不明。注意拼写：limb 末字母 b 不发音（同 \textbf{comb} 梳子、\textbf{dumb} 哑的、\textbf{climb} 爬）。
\end{origin}

\vw{rubber}{/ˈrʌbə(r)/}{n.}

\sense{n.}{橡胶}
\col{rubber boots 雨靴}
\ex{They wore rubber boots in the rain.}
\col{a rubber band 橡皮筋}

\sense{n.}{橡皮}
\ex{The pupil asked to borrow my rubber.}

\col{rubber（英）橡皮；（美）eraser}

\begin{origin}
rub（摩擦）+ -er——"用来摩擦的东西"。1770 年英国化学家普里斯特利发现这种物质能擦掉铅笔笔迹，遂命名 rubber（擦子）。天然橡胶本名 caoutchouc（来自南美原住民语，意为"会哭的树"）。
\end{origin}

\vw{conventional}{/kənˈvenʃənl/}{adj.}

\sense{adj.}{传统的，习惯的}
\col{conventional wisdom 传统观念}
\ex{Conventional wisdom says breakfast is the most important meal.}
\col{conventional methods 常规方法}

\sense{adj.}{依照惯例的，遵循习俗的}
\col{a conventional person 循规蹈矩的人}
\ex{His taste in art is fairly conventional.}

\col{conventional weapons 常规武器（区别于核武器）}

\begin{origin}
convention 来自拉丁语 \textit{conventio}（集会、约定），源于 \textit{convenire}（走到一起）：\textit{con-}（一起）+ \textit{venire}（来）。同源词：\textbf{convenient}（方便的）、\textbf{venue}（场地）、\textbf{event}（事件）。"大家聚到一起约定的做法"就是惯例。
\end{origin}

\vw{sensory}{/ˈsensəri/}{adj.}

\sense{adj.}{感觉的，感官的}
\col{sensory organs 感觉器官}
\ex{The eyes and ears are sensory organs.}
\col{sensory experience 感官体验}
\col{sensory overload 感官过载}

\begin{origin}
来自拉丁语 \textit{sensus}（感觉），源于 \textit{sentire}（感觉）。同源词：\textbf{sense}、\textbf{sensation}（轰动）、\textbf{sensible}（明智的）、\textbf{sensitive}（敏感的）。“五感”即 the five senses。
\end{origin}

\vw{patent}{/ˈpætnt; ˈpeɪtnt/}{n. \& adj.}

\sense{n.}{专利，专利证书}
\col{take out/obtain a patent on sth 申请专利}
\ex{He took out a patent on his invention.}
\col{hold a patent 持有专利}

\sense{adj. /ˈpeɪtnt/}{有专利的，受专利保护的}
\ex{a patented process 受专利保护的工艺}

\begin{origin}
来自拉丁语 \textit{patere}（张开、公开）$\to$ 法语 \textit{patente}。英国王室授特权时颁发"公开的信函"（letters patent），即不封口的公开文书——patent 由此得名"专利"。同源词：\textbf{patent} 与 "open" 同根（pate- 张开）。
\end{origin}

\vw{fuel}{/ˈfjuːəl/}{vt. \& n.}

\sense{n.}{燃料}
\col{fossil fuels 化石燃料}
\ex{Coal and oil are fossil fuels.}
\col{the fuel tank 油箱}

\sense{vt.}{增加，刺激}
\col{fuel inflation/speculation}
\ex{Rising wages fuel inflation.}
\col{fuel a debate/rumour 助长争论；加剧谣言}
\ex{The report fuelled speculation about a merger.}

\sense{vt.}{给……提供燃料}
\ex{The generator is fuelled by diesel.}

\begin{gram}[名词转动词（词性转换）]
\textit{fuel} 由名词"燃料"直接转作动词"添燃料、助长"（\textit{fuel a debate}），这种"零派生"构词叫 conversion。同类高频词：\textit{bottle}（装瓶）、\textit{hand}（递交）、\textit{head}（率领）、\textit{chair}（主持）、\textit{e-mail}（发邮件）。判断词性须看语境。
\end{gram}

\begin{origin}
来自拉丁语 \textit{focus}（壁炉）$\to$ 古法语 \textit{fouaille}（燃料）。同源词：\textbf{focus}（焦点）——拉丁语 focus 本义"壁炉"，因壁炉是房屋的"中心点"，引申为"焦点"。燃料与焦点，同出壁炉。
\end{origin}

\vw{disturb}{/dɪˈstɜːb/}{vt.}

\sense{vt.}{打扰}
\col{disturb sb}
\ex{Don't disturb me while I'm working.}
\col{Do not disturb. 请勿打扰}

\sense{vt.}{使不安}
\col{be disturbed by sth}
\ex{He was deeply disturbed by the news.}
\col{disturbing (adj.) 令人不安的}
\ex{a disturbing report}

\sense{vt.}{搅乱}
\col{disturb the peace 扰乱治安}
\ex{He was arrested for disturbing the peace.}

\begin{origin}
来自拉丁语 \textit{disturbare}（打散、扰乱）：\textit{dis-}（散）+ \textit{turbare}（搅乱），词根 \textit{turba}（人群、骚乱）。同源词：\textbf{turbine}（涡轮：旋转的机器）、\textbf{turbulent}（狂暴的）、\textbf{trouble}（麻烦：turbula 小骚乱）。"把平静搅乱"就是 disturb。
\end{origin}

\vw{tackle}{/ˈtækl/}{vt.}

\sense{vt.}{应付，处理，解决}
\col{tackle a problem/task}
\ex{The government must tackle unemployment.}
\ex{We need to tackle the problem at its root.}

\sense{vt.}{（体育中）铲球，擒抱}
\ex{He tackled the striker and won the ball.}

\begin{origin}
来自中古荷兰语 \textit{takel}（船上的索具）。由"抓取滑轮的绳索"引申出"抓住、着手处理"。水手用语进入日常英语的又一例（同 \textbf{deck} 甲板、\textbf{skipper} 船长）。
\end{origin}

\vw{head-on}{/ˌhed ˈɒn/}{adv. \& adj.}

\sense{adv.}{积极地，果断地}
\col{meet/confront sth head-on 迎难而上}
\ex{She decided to confront the problem head-on.}

\sense{adv.}{正面相撞地}
\ex{The two cars crashed head-on.}

\sense{adj.}{迎头相撞的；正面反对的}
\col{a head-on collision 迎头相撞}
\ex{Two drivers were killed in a head-on collision.}

\begin{origin}
head（头）+ on（对着）——"头对头"。与 \textbf{headlong}（头朝下地、鲁莽地）、\textbf{headfirst}（头朝前地）同类构词。
\end{origin}

\vp{go to great lengths to do sth}

\sense{v.}{竭尽全力做某事，不遗余力做某事}
\ex{She went to great lengths to make us feel welcome.}
\col{go to any lengths 不惜一切代价}
\ex{He would go to any lengths to get what he wants.}

\begin{origin}
length 是"长度"——"走到很长的距离"即费尽周折。同类表达：\textit{go out of one's way}（不辞辛劳）、\textit{spare no effort}（不遗余力）。
\end{origin}

\vw{adjust}{/əˈdʒʌst/}{vt. \& vi.}

\sense{vt.}{调整，调节}
\col{adjust sth}
\ex{You can adjust the height of the seat.}

\sense{vt. \& vi.}{适应，习惯}
\col{adjust to sth}
\ex{It took him a while to adjust to the new environment.}
\col{adjust oneself to sth}
\ex{She quickly adjusted herself to the new school.}

\sense{vt.}{整理}
\col{adjust one's clothes/hat}
\ex{He adjusted his tie before the interview.}

\col{adjustment (n.) 调整}

\begin{origin}
来自古法语 \textit{ajuster}（使接近恰当），源于 \textit{a-}（向）+ \textit{juste}（恰好的），再追溯到拉丁语 \textit{iustus}（公正的），词根 \textit{jus}（法）。同源词：\textbf{just}（恰好；公正的）、\textbf{justice}（正义）、\textbf{justify}（证明正当）。"调到恰好的位置"就是调整、适应。
\end{origin}

\vw{needle}{/ˈniːdl/}{n.}

\sense{n.}{针；注射针}
\col{a sewing needle 缝衣针}
\col{pine needles 松针}
\ex{The doctor gave me an injection with a needle.}

\sense{n.}{指针}
\col{the needle of a compass 罗盘指针}
\ex{The compass needle points north.}

\begin{origin}
源自古英语 \textit{nædl}，日耳曼语，词根可追溯到原始印欧语 *\textit{sneh-}（缝）。同源词：\textbf{sew}（缝纫）。"缝纫的针"就是 needle——needle 与 sew 同源。
\end{origin}

\vw{evidence}{/ˈevɪdəns/}{n. \& vt.}

\sense{n.}{证据，证明；物证，人证}
\col{evidence of/for sth}
\ex{There is no evidence of any crime.}
\col{give evidence 作证}
\ex{She gave evidence at the trial.}
\col{on the evidence of sth 根据……}
\ex{On the evidence of this report, sales are rising.}

\sense{vt.}{证明，表明}
\ex{His later career evidences his early talent.}

\begin{gram}[evidence 不可数]
\textit{evidence} 是不可数名词：\textcolor{WildStrawberry!80!black}{\ding{55}} \textit{an evidence}；应说 \textit{a piece of evidence}（一条证据）。谓语用单数：\textit{The evidence suggests ...}
\end{gram}

\begin{origin}
来自拉丁语 \textit{evidentia}，源于 \textit{evidens}（明显的）：\textit{e-}（向外）+ \textit{videre}（看见）。同源词：\textbf{video}（看见）、\textbf{visible}（可见的）、\textbf{provide}（预先看到 $\to$ 提供）。"看得见的东西"就是证据。
\end{origin}

\vp{date back to}

\sense{v.}{追溯到，始于，自……至今}
\ex{The tradition dates back to the Ming Dynasty.}
\col{date from 始于}
\ex{The castle dates from the 12th century.}
\col{date back + 时长}
\ex{The friendship dates back 30 years.}

\col{date back to 不用被动语态，一般用现在时}

\begin{origin}
date 由"写日期"引申为"确定年代"——"日期能追溯到……"。\textit{date back to} 强调从过去延续至今。
\end{origin}

\vw{solid}{/ˈsɒlɪd/}{adj.}

\sense{adj.}{固体的}
\col{solid food 固体食物}
\ex{The baby has started eating solid food.}

\sense{adj.}{实心的}
\ex{a solid ball of metal}

\sense{adj.}{结实的}
\col{a solid wall/building}
\ex{The house is built of solid brick.}

\sense{adj.}{纯色的}
\col{a solid colour 纯色}
\ex{She prefers solid colours to patterns.}

\col{solid evidence 确凿的证据（拓展）}
\ex{We have solid evidence against him.}

\begin{origin}
来自拉丁语 \textit{solidus}（坚实的、实心的）。同源词：\textbf{soldier}（士兵）——古罗马发军饷用的是 solidus 金币，"拿军饷的人"就是士兵；\textbf{consolidate}（巩固）。solid 与 soldier 同源。
\end{origin}

\vw{switch}{/swɪtʃ/}{vi. \& vt. \& n.}

\sense{vt.}{转变}
\col{switch from A to B}
\ex{She switched from history to economics.}
\col{switch roles/plans 换角色；改变计划}

\sense{vt.}{交换}
\col{switch seats 换座位}
\ex{Can we switch seats?}

\sense{vt.}{调班}
\col{switch shifts}
\ex{I asked to switch shifts with my colleague.}

\sense{n.}{开关，闸}
\col{a light switch 电灯开关}
\ex{The switch is by the door.}

\sense{n.}{转变，改变}
\col{make a switch}
\ex{The company made a switch to renewable energy.}

\col{switch sth on/off 打开/关闭}
\ex{Switch off the lights before you leave.}

\begin{origin}
来自中古荷兰语 \textit{swijch}（树枝、鞭子）。铁路时代，"可转动的岔道装置"（转辙器）让列车变轨——switch 遂由"树枝"引申为"开关、转换"。电子开关与"转变"都来自铁道转辙器。
\end{origin}

\vw{stainless}{/ˈsteɪnləs/}{adj.}

\sense{adj.}{不生锈的}
\col{stainless steel 不锈钢}
\ex{The pot is made of stainless steel.}

\begin{origin}
stain（污渍）+ -less（无）。stain 来自古法语 \textit{desteindre}（褪色），源于拉丁语 \textit{tingere}（染色）。同源词：\textbf{stain}、\textbf{tinge}（淡染）。"没有污渍、不生锈"就是 stainless。
\end{origin}

\vw{steel}{/stiːl/}{n.}

\sense{n.}{钢；钢铁工业}
\ex{The bridge is made of steel.}
\col{the steel industry 钢铁工业}
\ex{The steel industry employs thousands of workers.}
\col{nerves of steel 钢铁般的神经}
\ex{You need nerves of steel to be a surgeon.}

\col{steel oneself for sth 使自己坚强面对（拓展）}
\ex{She steeled herself for the difficult conversation.}

\begin{origin}
源自古英语 \textit{stīele}，日耳曼语（同源德语 \textit{Stahl}）。词源不确定，一说与"僵直"（stiff）相关——钢是"坚硬挺直"的金属。
\end{origin}

\vw{examine}{/ɪɡˈzæmɪn/}{vt.}

\sense{vt.}{检查，检验}
\col{examine sth closely}
\ex{The doctor examined the patient carefully.}

\sense{vt.}{审查，调查}
\col{examine the evidence}
\ex{The committee examined the evidence in detail.}

\sense{vt.}{考，测试}
\col{examine sb on sth}
\ex{Students will be examined on all the topics covered.}

\sense{vt.}{审问}
\col{examine a witness 询问证人}

\col{examination (n.) 考试；检查}

\begin{origin}
来自拉丁语 \textit{examinare}（称量、检验），源于 \textit{examen}（天平指针），词根 \textit{exigere}（衡量）：\textit{ex-}（出）+ \textit{agere}（驱动）。同源词：\textbf{exact}（精确的：称量出来的）。"考试"即"称量学问"。
\end{origin}

\vw{tongue}{/tʌŋ/}{n.}

\sense{n.}{舌头}
\col{Bite your tongue! 别乱说！}
\ex{The answer was on the tip of my tongue. 话到嘴边}
\ex{She was tongue-tied with shyness. 张口结舌}

\sense{n.}{语言}
\col{mother tongue 母语}
\ex{Her mother tongue is Spanish.}
\col{a foreign tongue 外语（书面）}

\begin{origin}
源自古英语 \textit{tunge}，日耳曼语（同源德语 \textit{Zunge}）。原始印欧语 *\textit{dnghu-}（舌头）在拉丁语中变成 \textit{lingua}（d 变 l）——同源词：\textbf{language}（语言）、\textbf{linguist}（语言学家）。tongue 与 language 同出一根："语言"就是"舌头的功夫"。
\end{origin}

\vw{pulse}{/pʌls/}{n. \& vi.}

\sense{n.}{脉搏}
\col{take/feel one's pulse 把脉}
\ex{The nurse took my pulse.}
\col{a strong/weak pulse}

\sense{vi.}{搏动，跳动}
\ex{Blood pulsed through his veins.}

\sense{vi.}{洋溢着}
\col{be pulsing with sth}
\ex{The city was pulsing with energy.}

\begin{origin}
来自拉丁语 \textit{pulsus}（推动），源于 \textit{pellere}（推、击），其反复动词 \textit{pulsare}（连连推）正是 \textbf{push}（推）的来源。同源词：\textbf{push}、\textbf{compel}（强迫）、\textbf{expel}（驱逐）。"心脏推动血液"的节奏就是脉搏。
\end{origin}

\vw{wrist}{/rɪst/}{n.}

\sense{n.}{手腕，腕关节}
\col{a wrist watch 手表}
\ex{He wore a silver wrist watch.}
\col{the wrist joint 腕关节}

\begin{origin}
源自古英语 \textit{wrist}，源于 \textit{wræstan}（扭）——"能扭动的关节"。同源词：\textbf{wrestle}（摔跤）、\textbf{writhe}（扭动）。手腕与摔跤，共享"扭"的意象。
\end{origin}

\vw{organ}{/ˈɔːɡən/}{n.}

\sense{n.}{器官}
\col{internal organs 内脏}
\ex{The heart is a vital organ.}
\col{organ donation 器官捐献}

\sense{n.}{风琴}
\col{a pipe organ 管风琴}
\ex{He played the organ in church.}

\sense{n.}{机构，机关}
\col{government organs 政府机关}

\begin{origin}
来自希腊语 \textit{organon}（工具、器械），源于 \textit{ergein}（工作），词根可追溯到原始印欧语 *\textit{werg-}（工作）——与英语 \textbf{work} 同源。身体的"工具"是器官，发声的"工具"是风琴，办事的"工具"是机关。
\end{origin}

\vw{severe}{/sɪˈvɪə(r)/}{adj.}

\sense{adj.}{十分严重的，极为恶劣的}
\col{severe weather/storms}
\ex{The region was hit by severe storms.}
\col{a severe injury/illness}
\ex{He suffered severe injuries in the crash.}

\sense{adj.}{严厉的}
\col{severe punishment/criticism}
\ex{The offenders face severe punishment.}

\sense{adj.}{艰难的}
\col{severe competition 激烈的竞争}
\ex{Severe competition has driven prices down.}

\col{severe pain 剧痛}

\begin{origin}
来自拉丁语 \textit{severus}（严格的、严肃的），词源不确定，可能与"冷硬"的感觉有关。同源词：\textbf{severity}（严重性）。
\end{origin}

\vw{anxiety}{/æŋˈzaɪəti/}{n.}

\sense{n.}{焦虑；担心}
\col{anxiety about/over sth}
\ex{She felt deep anxiety about the future.}
\col{a source of anxiety 焦虑的来源}
\ex{Money was a constant source of anxiety for him.}

\sense{n.}{渴望}
\col{anxiety to do sth}
\ex{his anxiety to please his parents}

\begin{origin}
来自拉丁语 \textit{anxietas}，源于 \textit{anxius}（焦虑的），词根 \textit{angere}（勒紧、使窒息），原始印欧语 *\textit{angh-}（紧）。同源词：\textbf{anger}（愤怒：内心被勒紧）、\textbf{anguish}（痛苦）、\textbf{angina}（心绞痛）。"心里发紧"就是焦虑——anxiety 与 anger 同源。
\end{origin}

\vw{abuse}{/əˈbjuːz/ \ vt. \ \ /əˈbjuːs/ \ n.}{vt. \& n.}

\sense{vt. /əˈbjuːz/}{滥用}
\col{abuse one's power/position}
\ex{He abused his position for personal gain.}

\sense{vt.}{虐待}
\col{abuse a child/an animal}
\ex{The prisoner was abused by the guards.}

\sense{vt.}{辱骂}
\ex{The driver abused the cyclist.}

\sense{n. /əˈbjuːs/}{滥用；虐待；辱骂}
\col{alcohol/drug abuse 酗酒；吸毒}
\ex{Drug abuse is a serious social problem.}
\col{child abuse 虐待儿童}

\begin{gram}[名动异读：清浊辅音转换]
\textit{abuse} 动词尾读浊音 /z/（/\textbf{ˈbjuːz}/），名词尾读清音 /s/（/\textbf{ˈbjuːs}/）。同类"名清动浊"的词对：\textbf{use}（v. /juːz/；n. /juːs/）、\textbf{excuse}、\textbf{house}（v. /haʊz/；n. /haʊs/）、\textbf{refuse}（v. /rɪˈfjuːz/；n. /ˈrefjuːs/）。可据此判断词性。
\end{gram}

\begin{origin}
来自拉丁语 \textit{abuti}（用尽、滥用）：\textit{ab-}（偏离）+ \textit{uti}（使用）。同源词：\textbf{use}、\textbf{usage}、\textbf{utilize}（利用）。"用偏了、用过头了"就是滥用。
\end{origin}

\vw{drug}{/drʌɡ/}{n.}

\sense{n.}{药物}
\ex{The doctor prescribed a new drug.}

\sense{n.}{毒品}
\col{take drugs 吸毒}
\ex{He was arrested for selling drugs.}
\col{a drug addict 瘾君子}
\ex{Drug addicts need treatment, not punishment.}
\col{drug trafficking 贩毒}

\begin{origin}
来自古法语 \textit{drogue}（药物），词源不确定，一说借自荷兰语 \textit{droog}（干的）——荷兰商人把"干燥的草药"（drogue）当药物贩卖。由"干草药"到"药物"再到"毒品"，一词记录了一部贸易史。
\end{origin}

\vw{phenomenon}{/fəˈnɒmɪnən/}{n.}

\sense{n.}{(pl. phenomena) 现象}
\col{a natural/social phenomenon}
\ex{Lightning is a natural phenomenon.}
\col{the phenomenon of sth}
\ex{the phenomenon of globalization}

\col{a phenomenon 奇才，非凡的人（拓展）}
\ex{The child is a phenomenon at chess.}

\begin{gram}[希腊语复数：phenomenon–phenomena]
\textit{phenomenon} 的复数不是 \textit{phenomenons}，而是希腊语复数 \textit{phenomena}。同类"希腊/拉丁语外来复数"：\textit{criterion–criteria}（标准）、\textit{medium–media}（媒介）、\textit{datum–data}（数据）、\textit{curriculum–curricula}（课程）。注意 phenomena 是复数，谓语用复数：\textit{These phenomena are rare.}
\end{gram}

\begin{origin}
来自希腊语 \textit{phainomenon}（显现的事物），源于 \textit{phainein}（显现），词根可追溯到原始印欧语 *\textit{bha-}（发光）。同源词：\textbf{phantom}（幻影）、\textbf{fantasy}（幻想）、\textbf{phase}（阶段：显现的侧面）。"显现出来的东西"就是现象。
\end{origin}

\vw{substance}{/ˈsʌbstəns/}{n.}

\sense{n.}{物质}
\col{a chemical substance}
\ex{Carbon is a substance found in all living things.}

\sense{n.}{要点}
\col{the substance of sth}
\ex{The substance of his speech was about tax reform.}
\col{in substance 实质上}
\ex{In substance, the two plans are the same.}

\begin{origin}
来自拉丁语 \textit{substantia}（本质），源于 \textit{substare}（站在下面）：\textit{sub-}（下）+ \textit{stare}（站）。同源词：\textbf{stand}、\textbf{stance}。中世纪哲学家翻译希腊语 \textit{ousia}（本质）时造了 substantia——"支撑事物之下的东西"就是本质、物质。
\end{origin}

\vw{symbol}{/ˈsɪmbl/}{n.}

\sense{n.}{象征}
\col{a symbol of sth}
\ex{The dove is a symbol of peace.}

\sense{n.}{符号，代号}
\col{a chemical symbol 化学符号}
\ex{\textit{H$_2$O} is the chemical symbol for water.}

\begin{origin}
来自希腊语 \textit{symbolon}（信物、标记），源于 \textit{symballein}（放在一起）：\textit{syn-}（一起）+ \textit{ballein}（投掷）。古希腊人缔约时各执信物的一半，日后合对以验明身份——symbol 本义是"合对信物"。象征即"一物代表另一物"。
\end{origin}

\vw{infection}{/ɪnˈfekʃn/}{n.}

\sense{n.}{传染，感染}
\col{an infection 感染}
\ex{The wound developed a serious infection.}
\col{spread infection 传播感染}
\col{infect (v.) 传染}
\ex{The virus infected thousands of computers.}
\col{infectious (adj.) 传染性的；有感染力的}
\ex{Her laughter is infectious.}

\begin{origin}
来自拉丁语 \textit{inficere}（浸染、污染）：\textit{in-}（进）+ \textit{facere}（做）——"做进去、染进去"。同源词：\textbf{infect}、\textbf{effect}、\textbf{affect}。"感染"即异物"浸染"进身体。
\end{origin}

\vw{department}{/dɪˈpɑːtmənt/}{n.}

\sense{n.}{（医院的）科，部门，局，系}
\col{the sales department 销售部}
\ex{She works in the marketing department.}
\col{the English department 英语系}
\col{the emergency department 急诊科}
\ex{He was rushed to the emergency department.}
\col{a department store 百货商店}

\begin{origin}
来自拉丁语 \textit{departire}（分开）：\textit{de-}（分）+ \textit{partiri}（分割），词根 \textit{pars}（部分）。同源词：\textbf{part}、\textbf{depart}（离开）、\textbf{partition}（隔断）。"分出来的部分"就是部门、科系。
\end{origin}

% ---------------- 选收的未标 * 词条 ----------------

\vw{gene}{/dʒiːn/}{n.}

\sense{n.}{基因}
\col{a dominant/recessive gene 显性/隐性基因}
\ex{The gene for eye colour is passed from parents to children.}

\begin{origin}
1909 年由丹麦遗传学家约翰森（Wilhelm Johannsen）造出，来自希腊语 \textit{genos}（出生、种族），词根可追溯到原始印欧语 *\textit{gen-}（生育）。同源词：\textbf{generate}（产生）、\textbf{gender}（性别）、\textbf{genius}（天赋）、\textbf{generation}（一代人）。基因即"决定出生形态的东西"。
\end{origin}

\vw{game-changing}{/ˈgeɪm tʃeɪndʒɪŋ/}{adj.}

\sense{adj.}{改变游戏规则的（指彻底改变事态发展的）}
\col{a game-changing technology/invention}
\ex{The smartphone was a game-changing invention.}

\begin{origin}
复合词：game（比赛）+ changing（改变）——"改变比赛规则的"，比喻"颠覆性的、划时代的"。
\end{origin}

\vw{virus}{/ˈvaɪrəs/}{n.}

\sense{n.}{病毒；（计算机程序中的）病毒}
\col{a computer virus 计算机病毒}
\ex{The virus spread rapidly through the school.}
\ex{The email contained a computer virus.}

\col{viral (adj.) 病毒的；疯传的（拓展）}
\ex{The video went viral.}

\begin{origin}
来自拉丁语 \textit{virus}（毒液、黏液）。同源词：\textbf{virulent}（剧毒的）、\textbf{viral}。罗马人用 virus 指蛇毒、毒液——现代病毒学沿用此词。
\end{origin}

\vw{abnormal}{/æbˈnɔːml/}{adj.}

\sense{adj.}{反常的，畸形的}
\col{abnormal behaviour}
\ex{The test revealed abnormal cells.}

\begin{origin}
\textit{ab-}（偏离）+ normal（正常）。normal 来自拉丁语 \textit{norma}（木匠的直角尺）——"合尺的就是正常的"。同源词：\textbf{norm}（规范）、\textbf{enormous}（e- 出 + norma $\to$ 超出尺度的 $\to$ 巨大的）。abnormal 与 enormous 一正一反，都源于"尺"。
\end{origin}

\vw{weapon}{/ˈwepən/}{n.}

\sense{n.}{武器，手段；兵器}
\col{a nuclear weapon 核武器}
\col{weapons of mass destruction 大规模杀伤性武器}
\col{a weapon against sth 对付……的武器}
\ex{Education is a powerful weapon against poverty.}

\begin{origin}
源自古英语 \textit{wæpen}，日耳曼语（同源德语 \textit{Waffe}）。词根可追溯到原始印欧语 *\textit{weip-}（抛掷）——最初的武器就是"投掷之物"。
\end{origin}

\vw{finding}{/ˈfaɪndɪŋ/}{n.}

\sense{n.}{调查发现，调研结果}
\col{the findings of a study}
\ex{The findings were published in a medical journal.}
\col{key findings 主要结论}

\begin{origin}
find（发现）+ -ing 名词化。find 源自古英语 \textit{findan}（找到），与"走、遭遇"的词根有关。"找到的东西"就是发现、结论。
\end{origin}

\vw{outgoing}{/ˌaʊtˈɡəʊɪŋ/}{adj.}

\sense{adj.}{外向的}
\col{an outgoing person}
\ex{She's friendly and outgoing.}

\sense{adj.}{即将离任的（拓展）}
\col{the outgoing president}
\ex{The outgoing CEO handed over to her successor.}

\begin{origin}
out（向外）+ going（走）——"向外走的"：性格上"向外走"即外向；职务上"向外走"即离任。
\end{origin}

\vw{grateful}{/ˈɡreɪtfl/}{adj.}

\sense{adj.}{感激的，表示感谢的}
\col{be grateful to sb for sth}
\ex{I'm grateful to you for your help.}
\col{grateful (adj.) 与 thankful 近义}

\begin{origin}
来自拉丁语 \textit{gratus}（令人愉悦的、感激的）——与 \textbf{grace}（优雅；恩典）同源。同源词：\textbf{gratitude}（感恩）、\textbf{congratulate}（祝贺）。“感激”即“记得别人施予的恩惠”。
\end{origin}

\vw{acupuncture}{/ˈækjupʌŋktʃə(r)/}{n.}

\sense{n.}{针刺疗法}
\col{acupuncture treatment 针灸治疗}
\ex{Acupuncture has been used in China for thousands of years.}

\begin{origin}
拉丁语 \textit{acus}（针）+ \textit{punctura}（刺）——"针+刺"的合成词。词根 \textit{pungere}（刺）的同源词：\textbf{acute}（尖锐的）、\textbf{pungent}（刺鼻的）、\textbf{puncture}（刺孔）。
\end{origin}

\vw{clinic}{/ˈklɪnɪk/}{n.}

\sense{n.}{私人诊所；门诊部；门诊时间}
\col{a dental clinic 牙科诊所}
\ex{The clinic is open from 9 to 5.}
\col{the clinic's opening hours 门诊时间}

\begin{origin}
来自希腊语 \textit{klinikē}（床边的医术），源于 \textit{klinē}（床）——古代医生"到病人床边"行医。同源词：\textbf{clinical}（临床的）。
\end{origin}

\vw{widespread}{/ˈwaɪdspred/}{adj.}

\sense{adj.}{分布广的，普遍的，广泛的}
\col{widespread support/use}
\ex{The policy gained widespread support.}
\ex{The disease is widespread in tropical regions.}

\begin{origin}
wide（宽）+ spread（铺开）——"铺得很宽"。"散布面广"即普遍。
\end{origin}

\vw{delete}{/dɪˈliːt/}{vt.}

\sense{vt.}{删去，删除}
\col{delete a file/message}
\ex{He deleted the email by mistake.}
\col{delete sth from sth}
\ex{Her name was deleted from the list.}

\begin{origin}
来自拉丁语 \textit{delere}（摧毁、抹去）。同源词：\textbf{indelible}（抹不掉的：in- 不 + delere）。"删除"即"抹掉字迹"。
\end{origin}

\end{document}
